\documentclass[conference]{IEEEtran}
\IEEEoverridecommandlockouts
\usepackage{cite}
\usepackage{amsmath,amssymb,amscd,accents,amsthm,amsfonts}%数式用のパッケージ%
\usepackage{multicol}%ノーテーション用のパッケージ%
\usepackage{graphicx}%画像挿入用%
\usepackage{ascmac,fancybox}%枠囲み用のパッケージ%
\usepackage{latexsym}%特殊文字の出力に必要%
\usepackage{here}%配置に使えるやつ%
\usepackage{bm}%太字にできる
\usepackage{float}%好きな位置に図を出せる
\usepackage{color}%文字に色つける
\usepackage{comment}%注釈つける
\usepackage{algorithm}%疑似コード
\usepackage{algpseudocode}%疑似コード
\usepackage{bbm}
\usepackage{booktabs}
\usepackage{cases}
\usepackage{mathtools}
\usepackage{mathrsfs}
\usepackage{breqn}
\usepackage{dsfont}
\usepackage{enumitem}
\newcommand{\argmax}{\mathop{\rm arg~max}}%argmaxを表示
\newcommand{\argmin}{\mathop{\rm arg~min}}%argminを表示
\renewcommand{\algorithmicrequire}{\textbf{Input:}}
\renewcommand{\algorithmicensure}{\textbf{Output:}}
\theoremstyle{definition}
\newtheorem{defi}{Definition}
\newtheorem{example}{Example}
\newtheorem{lem}{Lemma}
\newtheorem{result}{Result}
\newtheorem{thm}{Theorem}
\newtheorem{prop}{Proposition}
\newtheorem{caution}{Caution}
\newtheorem{assumption}{Assumption}
\newtheorem{corollary}{Corollary}
\newtheorem{remark}{Remark}
\def\BibTeX{{\rm B\kern-.05em{\sc i\kern-.025em b}\kern-.08em
    T\kern-.1667em\lower.7ex\hbox{E}\kern-.125emX}}
\hyphenpenalty=3000\relax
\exhyphenpenalty=3000\relax

\begin{document}

\title{A Stochastic Image Model Generating Object Regions and Approximate Algorithms of Bayes Optimal Semantic Segmentation
\thanks{This work was supported by JSPS KAKENHI, Grant Numbers JP23K03863.}
}

\author{\IEEEauthorblockN{Yilun Wang}
\IEEEauthorblockA{
\textit{SHOWA IKA KOHGYO Co. Ltd.}\\
Tokyo, Japan \\
yilun.wang@showaika.co.jp}
\and
\IEEEauthorblockN{Koshi Shimada}
\IEEEauthorblockA{
\textit{Waseda University}\\
Tokyo, Japan \\
i\_am\_koshi@suou.waseda.jp}
\and
\IEEEauthorblockN{Toshiyasu Matsushima}
\IEEEauthorblockA{
\textit{Waseda University}\\
Tokyo, Japan \\
toshimat@waseda.ac.jp}
}


\maketitle

\begin{abstract}

Stochastic model-based image segmentation offers statistical optimality under assumed priors, yet conventional models often lack a consistent generative process for object regions, labels, and pixel values. 
We propose a hierarchical model that defines region formation through a two-stage stochastic process. 
After generating quadtree-based blocks of varying sizes, the core novelty of our approach lies in constructing object regions through their stochastic connection, mediated by two intermediate graphs.
To circumvent the combinatorial explosion of parameters, we extend the concept of node-branching probabilities to connection probabilities for block combinations. 
This parametric approach enables efficient prior assignment for all possible region patterns with a limited parameter set. 
While we successfully derive the posterior distribution, calculating the exact Bayes-optimal solution is computationally intensive. 
We propose an iterative approximation algorithm for posterior probability updates. Numerical experiments and real-world applications verify the effectiveness and optimality of our method.

\end{abstract}

\section{Introduction}

Image segmentation has been an area of active research in the field of computer vision. 
Specifically, it covers a wide range of tasks, from dividing an image into homogeneous regions to semantic segmentation and instance segmentation, which partition an image by object class~\cite{Minaee_survey}. 

Previous methods for identifying object regions within an image have adopted several approaches. 
First, there are discriminative and algorithmic approaches including signal processing, clustering methods, and neural network models. 
For example, many of these frameworks partition the image into segments \cite{Uijlings_SelectiveSearch} and then identify labels for each region 
~\cite{Tsai_fourier, Long_CNN}. 
While these methods are effective in practice, they do not assume an underlying stochastic model for the data; consequently, their statistical optimality remains theoretically unproven.
Second, there are hybrid approaches that partially incorporate stochastic models into algorithmic frameworks.
For instance, some methods perform initial label estimation under pixel-wise stochastic models, followed by signal processing to determine the object regions~\cite{Carson_Blobworld, Zhang_MRF}. 
Conversely, other methods involve initially partitioning the image into small regions using discriminative algorithms and then estimating object regions by merging or splitting these small regions based on models of their connectivity or label correlations~\cite{Yao_OMRF, Zheng_OMRF, ghosh_ddCRP}. 

However, research on a consistent stochastic model that treats not only labels and pixel values but also the regions themselves as random variables is limited, although unified modeling of the entire generation process is essential for guaranteeing statistical optimality.
Our previous research~\cite{wang_isita2024} assumed a consistent stochastic model where square regions were constructed via a quadtree, with labels and pixel values generated independently for each region. 
By setting parametric prior distributions, we formulated Bayes-optimal segmentation and derived an algorithm that can be computed analytically while maintaining Bayesian optimality. 
However, while this succeeded in modeling square regions, it suffered from limited expressive power, making it difficult to represent complex and diverse object shapes.

To address this limitation, this study introduces a new consistent stochastic model capable of representing regions with arbitrary and diverse shapes. 
Although various methods for constructing regions are conceivable, we determine regions through the following two-step process: 
In the first step, as in our previous research~\cite{wang_isita2024}, we determine square regions of various sizes using a quadtree. 
In the second step, we newly propose a model that constructs regions by combining square regions based on two graph structures. 
By assuming a model in which square regions formed by a quadtree and the regions constructed by their combination occur according to certain probabilities, we can establish a consistent stochastic model for the image that represents the generation of regions with various shapes.

However, assigning prior probabilities individually to all possible patterns of quadtrees and their combined regions would result in the number of parameters growing exponentially with respect to the image size. 
In our previous research~\cite{wang_isita2024}, by setting branching probabilities as parameters for each node of the quadtree, we were able to assign prior probabilities to all possible quadtrees using only a small number of parameters. 
Therefore, in this study, we similarly propose assigning prior probabilities to all region shapes with a small number of parameters by assigning connection probabilities to each vertex of the graph.

Based on the proposed model above, we present a Bayes-optimal solution for semantic segmentation. 
By assuming parametric prior distributions, we were also able to derive the posterior distribution. 
However, since the computational cost of the Bayes-optimal solution remains high and exact computation is difficult, we propose two types of efficient approximation algorithms. 
Through numerical experiments using synthetic data, we evaluate the convergence of these approximation methods to the Bayes-optimal solution, and further demonstrate the practicality of the proposed model through application to real datasets.

\begin{comment}
The structure of this paper is as follows. 
Section II describes the proposed stochastic model for an image. 
Section III presents the optimal semantic segmentation solution under the assumed model. 
Sections IV and V define the prior distributions for the quadtree and regions, and describe efficient methods for calculating their posterior distributions. 
Section VI proposes an approximation algorithm for calculating the Bayes-optimal solution. 
Finally, in Section VII, we demonstrate the effectiveness and optimality of the proposed algorithm through experimental results.
\end{comment}

\section{The Proposed Stochastic Model}
First, we state the basic notation for images.

\begin{defi}
A two-dimensional image is composed of pixels arranged in a grid of size $D\times D$, where $D = 2^{d_{\mathrm{max}}}$ and $d_{\mathrm{max}} \in \mathbb{Z}_{\geq 0}$. 
Let $(i,j)$ denote pixel coordinates, and define the set of all pixels as $s_{\text{root}} = \{(i,j) \in \mathbb{Z}_{\geq 0}^2 \mid 0 \leq i, j \leq D-1\}$. 
\end{defi}
\begin{defi}
Let $\mathcal{Y}$ be a set of possible pixel values, with $y_{(i,j)} \in \mathcal{Y}$ denoting the value at $(i, j)$. 
The image is represented by the matrix $Y = (y_{(i,j)})_{0 \le i, j \le D-1}$.
\end{defi}

We assume a stochastic model where each pixel has a true label indicating its object class within the image $Y$. 
Each label is associated with an object region composed of square blocks of various sizes. 
Since the true labels and regions are unobserved, they are assumed to follow a certain prior distribution.


\subsection{The Region Model}

In this section, to represent regions with diverse shapes, we assume a mathematical model that constructs regions in two stages: while the first stage generates square blocks of various sizes via a quadtree following the previous work\cite{wang_isita2024}, the second stage introduces our newly proposed approach to combine them.

\textbf{[Stage 1]} We first describe the model for partitioning the image into square regions via a quadtree~\cite{wang_isita2024, nakahara_quadtree}.


\begin{defi}
Let $d$ be an integer satisfying $0 \leq d \leq d_{\mathrm{max}}$, and let $x_d, y_d \in \{0,1\}$. A subset $s_{(x_1,y_1)...(x_d,y_d)}$ of $s_{\text{root}}$ is defined as follows:
\begin{align}
&s_{(x_1,y_1)...(x_d,y_d)}\coloneqq \Biggl\{(i,j)\in\mathbb{Z}^2 \Biggm| \sum_{d'=1}^d \frac{x_{d'}}{2^{d'}}\leq \frac{i}{2^{d_{\mathrm{max}}}}<\nonumber\\
&\Biggl(\sum_{d'=1}^d \frac{x_{d'}}{2^{d'}} + \frac{1}{2^d}\Biggr),\sum_{d'=1}^d \frac{y_{d'}}{2^{d'}}\leq \frac{j}{2^{d_{\mathrm{max}}}}<\left(\sum_{d'=1}^d \frac{y_{d'}}{2^{d'}} + \frac{1}{2^d}\right)\Biggr\}. 
\end{align}
 Furthermore, the set of all possible blocks is defined as:
\begin{eqnarray}
\mathcal{S}_{d_{\mathrm{max}}} \coloneqq \{s_{\text{root}}, s_{(00)}, ..., s_{(00)(00)}, ..., s_{(11)(11)...(11)}\}.
\end{eqnarray}
\end{defi}

\begin{defi}
  Let $T$ be a rooted quadtree whose nodes are elements of $\mathcal{S}_{d_{\mathrm{max}}}$ (Fig.~\ref{quadtree}). 
  Let $\mathcal{L}_T$ and $\mathcal{I}_T$ denote the sets of leaf and internal nodes in $T$, respectively. 
  The depth of a node $s$ is denoted by $\mathrm{depth}(s)$, with $\mathrm{depth}(s_{\text{root}})=0$. 
  Let $\mathcal{T}_{d_{\text{max}}}$ be the set of all possible quadtrees with a maximum depth $d_{\text{max}}$ (Fig.~\ref{quadtree_set}). 
  The set of child nodes of a node $s$ is denoted by $\mathrm{Ch}(s)$.
\end{defi}

\begin{figure}[tb]
  \centering
  \includegraphics[width=6.5cm]{quadtree.png}
  \caption{The correspondence between the quadtree and the square blocks} \label{quadtree}
\end{figure}

\begin{figure}[tb]
  \centering
  \includegraphics[width=8.5cm]{Set_of_Quadtree.png}
  \caption{The set of all quadtrees with depth up to $d_{\text{max}}$} \label{quadtree_set}
\end{figure}


Since the quadtree $T\in\mathcal{T}_{d_{\text{max}}}$ corresponding to the image $Y$ is unknown, we treat it as a random variable following a prior distribution.
However, assigning a prior probability to each $T \in \mathcal{T}_{d_{\text{max}}}$ individually would result in the number of parameters growing exponentially with the image size. 
To address this, we assume a prior distribution where the quadtree is generated recursively from the root node $s_{\text{root}}$, using a branching probability $g_s$ for each $s \in \mathcal{S}_{d_{\text{max}}}$ as a parameter~\cite{wang_isita2024, nakahara_quadtree}:

\begin{assumption}
  For all nodes $s\in \mathcal{S}_{d_{\text{max}}}$, we set a parameter $g_s\in[0,1]$. 
  Then, the prior distribution of the quadtree $T$ is expressed by the following equation:
  \begin{align}
    p(T)=\prod_{s\in \mathcal{L}_T}(1-g_s)\prod_{s'\in \mathcal{I}_T}g_{s'}.
  \end{align}
  Here, $g_s$ represents the branching probability at node $s$. We set $g_s=0$ when $\mathrm{depth}(s)=d_{\mathrm{max}}$.
  \label{prior_dist_T}
\end{assumption}

This formulation successfully assigns prior probabilities to all quadtrees in $\mathcal{T}_{d_{\text{max}}}$, and constitutes a valid probability distribution \cite{nakahara_quadtree,nakahara_full_rooted_trees}.

\textbf{[Stage 2]} The object regions are constructed by combining the leaf nodes $s\in \mathcal{L}_T$ of a quadtree $T$ generated according to the prior distribution $p(T; \mathbf{g})$.

\begin{defi}
  A region $r$ is a set of pixels $(i, j)$ formed by the union of one or more leaf nodes $s \in \mathcal{L}_T$ of a quadtree $T$. 
  $R$ is defined as a set of regions $\{r_1, r_2, \dots, r_{|R|}\}$ that are mutually disjoint and whose union covers all pixels $s_{\text{root}}$. 
  The set of all possible regions that can be constructed from the leaf nodes of a quadtree $T$ is denoted by $\mathcal{R}_T$ (Fig.~\ref{region}). 
\end{defi}

\begin{figure}[tb]
  \centering
  \includegraphics[width=8.5cm]{region.png}
  \caption{Regions $R$ and the set $\mathcal{R}_T$ defined by combinations of leaf nodes in a quadtree $T$.} \label{region}
\end{figure}

Since the set of regions $R\in \mathcal{R}_T$ corresponding to the image $Y$ is unknown, we assume that $R$ is a random variabele.
However, directly assigning priors to all $R \in \mathcal{R}_T$ results in a combinatorial explosion of parameters. Instead, our proposed model introduces latent connection variables for leaf nodes $\mathcal{L}_T$, forming regions through a process involving two graphs.

\begin{defi}
   For a given quadtree $T$, let $G_T = (\mathcal{L}_{T}, E_T)$ be an undirected adjacency graph, where the edge set $E_T$ consists of all spatially adjacent pairs of leaf nodes. The set of adjacent nodes for $s \in \mathcal{L}_{T}$ is defined based on this graph as $\mathcal{E}_{T}(s) = \{s' \mid \{s, s'\} \in E_T\}$.
   \label{adjacent_graph}
\end{defi}

\begin{defi}
  Let $c_s \in \mathcal{E}_T(s) \cup \{s\}$ be a connection target selected by each node $s \in \mathcal{L}_{T}$ and the set of them be $\bm{c} = \{c_s\}_{s \in \mathcal{L}_T}$. A directed connection graph $G_T(\bm{c}) = (\mathcal{L}_{T}, E_T(\bm{c}))$ is formed with edges $E_T(\bm{c}) = \{(s, c_s) \mid s \in \mathcal{L}_{T}\}$, where each node has an out-degree of 1. 
  Each connected component of $G_T(\bm{c})$ defines a single region $r$. The set of regions formed by the connections $\bm{c}$ is denoted as $R=R(\bm{c})$ (Fig.~\ref{connect}).
  \label{connection_graph}
\end{defi}
In this study, we assume that the connection target $c_s$ for each leaf node is an independent random variable determined collectively and follows the prior distribution below.

\begin{assumption}
  Each connection target $c_s$ of a leaf node $s\in\mathcal{L}_T$ is decided according to the following probability:
  \begin{align}
  p(c_s \mid T) = 
  \begin{cases} 
  \frac{w(s, c_s)}{\alpha + \sum_{s'\in\mathcal{E}_T(s)}w(s, s')}, & c_s \in \mathcal{E}_T(s), \\ 
  \frac{\alpha}{\alpha + \sum_{s'\in\mathcal{E}_T(s)}w(s, s')}, & c_s = s.
  \end{cases}
  \end{align}
  Here, the function $w(s, c_s)$ is an asymmetric function that returns the connection weight for the adjacent leaf node $c_s \in \mathcal{E}_T(s)$ of leaf node $s$, and the constant $\alpha$ is the weight for selecting itself as the connection target.
The probability of generating the set of regions $R(\bm{c})$ given a quadtree $T$ coincides with the joint probability of the connection targets $\bm{c}$ and is calculated as follows:
\begin{align}
p(R(\bm{c}) \mid T) = p(\bm{c} \mid T) = \prod_{s \in \mathcal{L}_T} p(c_s \mid T).
\end{align}
\label{prior_of_connection}
\end{assumption}


\begin{thm}
  For any set of regions $R \in \mathcal{R}_T$, there exists at least one configuration of connection targets $\bm{c}$ such that $R(\bm{c}) = R$. Furthermore, the probability distribution $p(R \mid T)$ satisfies the normalization condition on $\mathcal{R}_T$:
  \begin{align}
    \sum_{R \in \mathcal{R}_T} p(R \mid T) \coloneqq  \sum_{R \in \mathcal{R}_T} \sum_{\bm{c}\in \{\bm{c'}\mid R(\bm{c'}) = R\}} p(\bm{c} \mid T) = 1.
  \end{align}
\end{thm}

\begin{remark}
While Definitions~\ref{adjacent_graph} and~\ref{connection_graph} and Assumption~\ref{prior_of_connection} align with an aspect of the distance-dependent Chinese Restaurant Process (ddCRP) framework~\cite{ghosh_ddCRP}, our approach provides a more mathematically rigorous description through the use of explicit graph structures.
\end{remark}

\begin{figure}[tb]
  \centering
  \includegraphics[width=8cm]{connect.png}
  \caption{Adjacency graph $G_T$ constructed from a given quadtree $T$, connection graph $G_T(\mathbf{c})$ for a specific connect target $\mathbf{c}$, and the set of regions $R = R(\mathbf{c})$ formed by their respective connected components.}
  \label{connect}
\end{figure}

\subsection{The Label Model}

\begin{defi}
  Let a set of label values be $\mathcal{X}$ and $x_{(i,j)} \in \mathcal{X}$ be a label value at a pixel $(i,j)$. 
  A label matrix is denoted as $X=(x_{(i,j)})_{0\leq i,j \leq D-1} \in \mathcal{X}^{D \times D}$.
  In this paper, pixels within the same region $r \in R(\bm{c})$ are assigned an identical label $x_r \in \mathcal{X}$. 
  This implies $x_{(i,j)} \coloneqq x_r$ for any $(i,j) \in r$.
\end{defi}

\begin{assumption}
  The label matrix $X\in\mathcal{X}^{D \times D}$ is assumed to be a random variable following a prior distribution $p(X \mid R(\bm{c}))$. 
  We assume that each label $x_r$ is assigned independently for each region $r \in R(\bm{c})$ as follows:
  \begin{align}
    p(X \mid R(\bm{c})) = \prod_{r \in R} p(x_r \mid R(\bm{c})). \label{prior_X_1}
  \end{align}
  We further assume that $p(x_r \mid R)$ depends solely on the features of region $r$. 
  \begin{align}
    \forall R \in \mathcal{R}_T, \forall r \in R, \quad p(x_r \mid R) = p(x_r). \label{prior_X_2}
  \end{align}
  Based on equations (\ref{prior_X_1}) and (\ref{prior_X_2}), the prior distribution is given by
  \begin{align}
    p(X \mid R(\bm{c})) = \prod_{r \in R(\bm{c})} p(x_r).
  \end{align}
\end{assumption}

When shape features of regions differ for each object, this prior knowledge can be incorporated into the prior distribution of labels~\cite{Dryden_ShapePrior, Gould_CRF}.


\subsection{The Pixel Value Model}


\begin{assumption}
Let $\bm{y}_r$ be the vector of pixel values in region $r$, where the elements are arranged in raster-scan order. Given the set of regions $R(\bm{c})$ and the label matrix $X$, the pixel value matrix $Y\in\mathcal{Y}^{D\times D}$ is assumed to be a random variable following a probability distribution $p(Y\mid X, R(\bm{c}))$. We assume that each pixel value vector $\bm{y}_r$ is generated independently for each region $r\in R(\bm{c})$ as follows:
\begin{align}
  p\left(Y \mid X, R(\bm{c})\right) = \prod_{r \in R(\bm{c})} p(\bm{y}_r \mid x_r).
\end{align}
\end{assumption}

When color or texture features of regions differ for each object, this prior knowledge can be incorporated into the probability distribution of pixel values~\cite{Carson_Blobworld, Oe_2DAR}. 





\section{Bayes-Optimal Solution for Semantic Segmentation}

Based on Bayes decision theory~\cite{berger_bayes_decision_theory}, we estimate the labels $X$ from the observed image $Y$.
We define the loss function for the estimated labels $\hat{X}$ as the average value of the 0-1 loss at each pixel:
\begin{align}
  \text{Loss} = \frac{1}{D \times D} \sum_{0 \leq i, j \leq D-1} \mathds{1} \{\hat{x}_{(i, j)} \neq x_{(i, j)}\}. \label{loss_function}
\end{align}
In this case, the Bayes-optimal label estimation is obtained by selecting the labels that maximize the following marginal posterior distribution:
\begin{align}
  &\hat{x}_{(i,j)} = \argmax_{x_{(i, j)} \in \mathcal{X}} \Biggl\{\sum_{T \in \mathcal{T}_{d_{\text{max}}}}  p(T \mid Y) \nonumber \\
  &\times \sum_{\bm{c} \in \prod\limits_{s \in \mathcal{L}_T} (\mathcal{E}_T(s) \cup \{s\})} p(\bm{c} \mid T, Y) p(x_{(i, j)} \mid R(\bm{c}), Y)\Biggr\}. \label{bayes_optimal_solution}
\end{align}


\section{Posterior Distribution of Regions}

\subsection{Posterior Distribution of a Quadtree}

Calculating the posterior $p(T \mid Y)$ requires marginalizing over labels $X$ and connection targets $\bm{c}$, which is analytically intractable. 
To address this, we approximate the distribution as follows.

The marginal likelihood $p(\bm{y}_{s}; g_s)$ for node $s$ is defined recursively, approximating the marginalization over $X$ and $\bm{c}$ using a non-informative prior:
\begin{align}
  &p(\bm{y}_{s}; g_s) \approx \nonumber \\
  &\begin{dcases}
    \phi(s), & \mathrm{depth}(s) = d_{\mathrm{max}}, \\
    (1-g_{s}) \phi(s) + g_{s} \prod_{s'\in \mathrm{Ch}(s)} p(\bm{y}_{s'}; g_{s'}), & \mathrm{depth}(s) < d_{\mathrm{max}}.
  \end{dcases}
\end{align}
where $\phi(s) \coloneqq \frac{1}{|\mathcal{X}|} \sum_{x_s \in \mathcal{X}} p(\bm{y}_{s} \mid x_{s})$ is the average likelihood over the label space.

\begin{thm}
  Given an image $Y$, the posterior distribution of the quadtree $T$ is given by the following~\cite{wang_isita2024, nakahara_quadtree}:
  \begin{align}
    p(T\mid Y)
    = \prod_{s\in\mathcal{L}_T}(1 - g_{s\mid Y})
      \;\prod_{s\in\mathcal{I}_T} g_{s\mid Y},
    \label{m's posterior}
  \end{align}
  where
  \begin{eqnarray}
    g_{s\mid Y}\coloneqq
    \begin{dcases}
      g_{s}, & \mathrm{depth}(s) = d_{\mathrm{max}},\\[6pt]
      \displaystyle\frac{g_{s}
        \displaystyle\prod_{s'\in \mathrm{Ch}(s)}p(\bm{y}_{s'}; g_{s'})}
        {p(\bm{y}_{s}; g_s)}, & \mathrm{depth}(s) < d_{\mathrm{max}}.
    \end{dcases}
  \end{eqnarray}
  \label{posterior_dist_T}
\end{thm}

\subsection{Posterior Distribution of Connection Targets}

Exact computation of $p(\bm{c} \mid T, Y)$ is costly.
However, since the conditional distribution for a single leaf node $s$ is tractable, we can approximate the posterior by iteratively updating $c_s$, as in Gibbs sampling or variational inference.

\begin{defi}
  Let $\bm{c}_{-s} \coloneqq \{c_{s'}\}_{s'\in\mathcal{L}_T\setminus \{s\}}$ be the set of connection targets excluding leaf node $s$. In addition,  $r_s(\bm{c})$ denote the region containing node $s$ under the connection $\bm{c}$.
\end{defi}

\begin{thm}
  Given $Y$, $T$, and $\bm{c}_{-s}$, the conditional probability of $c_s = s' \in \mathcal{E}_T(s)$ is determined by the likelihood ratio:
\begin{align}
  &p(c_s = s' \mid \bm{c}_{-s}, T, Y) \nonumber \\
  &\quad \propto \begin{dcases}
    \begin{array}{@{}l@{}}
      p(c_s = s' \mid T) \frac{p(\bm{y}_{r_s(\bm{c}_{-s}) \cup r_{s'}(\bm{c}_{-s})})}{p(\bm{y}_{r_s(\bm{c}_{-s})}) p(\bm{y}_{r_{s'}(\bm{c}_{-s})})} \\
      \multicolumn{1}{r}{, \, r_s(\bm{c}_{-s}) \neq r_{s'}(\bm{c}_{-s})},
    \end{array} \\
    p(c_s = s' \mid T) , \quad r_s(\bm{c}_{-s}) = r_{s'}(\bm{c}_{-s}).
  \end{dcases}
\end{align}
  Here, the marginal likelihood $p(\bm{y}_r)$ of the image data $\bm{y}_r$ in region $r$ can be calculated as follows:
  \begin{equation}
    p(\bm{y}_r) = \sum_{x_r \in \mathcal{X}} p(x_r) p(\bm{y}_r \mid x_r).
  \end{equation}
  \label{posterior_R}
\end{thm}

\section{Approximation Algorithms of Bayes Optimal Solution}

Despite parameter-efficient priors, the Bayes-optimal solution (\ref{bayes_optimal_solution}) remains exponentially complex relative to the image size. While exact MCMC is theoretically possible, its computational cost is often prohibitive. Consequently, we present two efficient approximation algorithms for posterior updates: Gibbs sampling and iterative MAP estimation.
Since both algorithms do not involve combinatorial operations, their computational complexity is of polynomial order relative to the image size.

\subsection{Gibbs Sampling Method}

\begin{enumerate}
  \item Calculate the quadtree $\hat{T}$ that maximizes the posterior probability~\cite{nakahara_full_rooted_trees}.
      \begin{align}
        &\max_{T \in \mathcal{T}_{d_{\mathrm{max}}}} p(T \mid Y) = \psi (s_{\text{root}}),\\
        &\psi(s) = \max \left\{ 1 - g_{s \mid Y}, \ g_{s \mid Y} \prod_{s' \in \mathrm{Ch}(s)} \psi(s') \right\}.
      \end{align}
\end{enumerate}

\begin{enumerate}[resume]
  \item For all leaf nodes $s \in \mathcal{L}(\hat{T})$ of $\hat{T}$, set the initial state of the connection target $c_s \in\mathcal{E}\cup \{s\}$. For example, set self-connections ($c_s = s$) as the initial state, where all nodes refer to themselves.
\end{enumerate}

\begin{enumerate}[resume]
  \item For each leaf node $s \in \mathcal{L}(\hat{T})$, sample the connection target based on $p(c_s \mid \bm{c}_{-s}, \hat{T}, Y)$ and update $\hat{\bm{c}}$ to a new connection target $c_s$. Convergence to the posterior distribution is sought by iteratively sampling for all leaf nodes in sequence. 
\end{enumerate}

\begin{enumerate}[resume]
  \item Estimate the optimal labels $\hat{x}_r$ to be assigned to each region $r \in R(\hat{\bm{c}})$: $\displaystyle \hat{x}_r = \argmax_{x_r \in \mathcal{X}} p(x_r \mid Y).$
\end{enumerate}

\label{gibbs_sampling}

\subsection{Updating Posterior Probability of Regions and Labels}

\begin{enumerate}
  \item Similarly to the previous section, calculate the quadtree $\hat{T}$ that maximizes the posterior probability and set the initial values for the connection variables.
  \item Iteratively update the connection targets $\hat{\bm{c}}=\{\hat{c}_s\}_{s\in\mathcal{L}_T}$ and labels $\hat{X}=\{\hat{x}_r\}_{r\in \mathcal{R}(\hat{\bm{c}})}$ as follows:
  \begin{itemize}
    \item For each $r\in R(\hat{\bm{c}})$, $\displaystyle \hat{x}_r = \argmax_{x_r \in \mathcal{X}} p(x_r \mid Y)$.
    \item For each $s\in\mathcal{L}_T$, $\displaystyle \hat{c}_s = \argmax_{c_s\in\mathcal{E}\cup \{s\}} p(c_s|\bm{c}_{-s}, \hat{T}, \hat{X}, Y)$.
  \end{itemize}
\end{enumerate}

\label{update_argmax}

\section{Experiments}

\subsection{Numerical Experiments}

We evaluate the proposed algorithms' approximation of the Bayes-optimal solution using synthetic data generated from our model. We used $N_{\text{train}}=50$ training and $N_{\text{test}}=10$ test images of size $128 \times 128$.

The parameter settings of the proposed model for generating synthetic data are as follows. 
First, the parameters $\bm{g}$ of the prior distribution for the quadtree $T$ were set to common values for each depth: 0.9, 0.8, 0.7, 0.6, 0.5, 0.4, 0.3, and 0, starting from depth 0. 
Then, for the prior distribution of region $R$, the self-connection weight was set to $\alpha=1.0\times 10^{-8}$. To favor larger nodes in region interiors and smaller nodes near boundaries, the connection weight function was set as $w(s, c_s) = \exp(\eta (\mathrm{depth}(s) - \mathrm{depth}(c_s)))$ with $\eta=10.0$. 
Moreover, the labels $\mathcal{X}=\{0,1,2\}$ are modeled via shape features $\bm{\phi}(r) = (\log \text{Area}(r), \text{Circularity}(r))$ for each region $r \in R$, assuming $p(x_r) \propto p(\bm{\phi}(r) \mid x_r)$ where each feature follows a label-dependent Gaussian distribution. 
The parameters were set such that regions with distinct combinations of size and circularity are assigned different labels.
Finally, the pixel value set is defined as $\mathcal{Y}=\{0,1,...,255\}^3$, and pixel values $\bm{y}_r$ are assumed to follow a pixel-wise independent normal distribution: $p(\bm{y}_r\mid x_r) = \prod_{(i,j)\in r} \mathcal{N}^3(y_{(i,j)};\boldsymbol{\mu}_{x_r}, \Sigma_{x_r})$. 
Pixel values corresponding to each label are generated with distinct variances and slightly differing intensities of red.

An example of the synthetic data generated with these settings is shown in Fig.~\ref{aritificial_data_exp_loss}(a). 
Parameters for the quadtree, labels, and pixels were estimated from $N_{\text{train}}$ training samples \cite{wang_isita2024}, while the region prior parameters were assumed to be known. 

Labels for the test data were then estimated using the proposed algorithms.
Fig.~\ref{aritificial_data_exp_loss}(b) shows the average test loss (\ref{loss_function}) over 10 updates converging toward zero, verifying optimal solution approximation. The algorithm in Sec.~\ref{update_argmax} converged faster, but Gibbs sampling (Sec.~\ref{gibbs_sampling}) theoretically yields higher accuracy with more samples.

\begin{figure}[tb]
  \centering
  \includegraphics[width=8cm]{aritificial_data_exp_loss.png}
  \caption{(a) Generated data by the proposed stochastic model, (b)Trend of average loss on test data by algorithms in Section \ref{gibbs_sampling} and Section \ref{update_argmax}}
  \label{aritificial_data_exp_loss}
\end{figure}

\label{exp_artificial_data}

\subsection{Application for Real Image}

\begin{figure}[tb]
  \centering
  \includegraphics[width=7.5cm]{ikonos_result.png}
  \caption{The result of segmentation for IKONOS} \label{ikonos_result}
\end{figure}

This section evaluates the proposed algorithm using real-world IKONOS satellite data \cite{IKONOS}. We assume that satellite imagery possesses a spatial structure consistent with our model, similar to the synthetic data described in Section \ref{exp_artificial_data}.

For this data, the set of labels is $\mathcal{X}=\{0,1\}$ and the set of pixel values is $\mathcal{Y}=\{0, 1, \dots, 255\}^3$.
Furthermore, we assumed the same settings as in Section \ref{exp_artificial_data} for the connection weight function $w(s, c_s)$ incorporated into the prior distribution for regions $R$, the prior distribution for labels $X$, and the likelihood function for pixel values $Y$.

We compared our method against GMM~\cite{Dempster_GMMEM}, MRF~\cite{Zhang_MRF}, Object-based MRF~\cite{Yao_OMRF} and the quadtree model ~\cite{wang_isita2024}. The estimation results are presented in Fig.~\ref{ikonos_result}.

As shown in Fig.~\ref{ikonos_result}, the proposed model suppresses mosaic-like noise and produces more natural segments than conventional models. This result suggests that explicit modeling of region shapes allows for superior segmentation, capturing complex structures that models without region-generation priors fail to represent.


\section{Conclusion}

  We proposed a consistent stochastic model capable of representing diverse region shapes. Regions are modeled as stochastic combinations of leaf nodes of a quadtree. To suppress the number of parameters, we introduced a prior distribution incorporating branching probabilities and connection probabilities for individual nodes. Finally, we derived the Bayes-optimal solution and its approximation algorithms, verifying their effectiveness through experiments.








\begin{thebibliography}{00}
  \bibitem{Minaee_survey} S. Minaee, Y. Boykov, F. Porikli, A. Plaza, N. Kehtarnavaz and D. Terzopoulos, ``Image Segmentation Using Deep Learning: A Survey,'' in IEEE Transactions on Pattern Analysis and Machine Intelligence, vol. 44, no. 7, pp. 3523-3542, 1 July 2022.
  \bibitem{Uijlings_SelectiveSearch} J. R. R. Uijlings, K. E. A. van de Sande, T. Gevers, and A. W. M. Smeulders, ``Selective Search for Object Recognition,'' Int. J. Comput. Vis., vol. 104, no. 2, pp. 154–171, Sep. 2013.
  \bibitem{Tsai_fourier} D.-M. Tsai and W.-L. Chen, ``Coffee plantation area recognition in satellite images using Fourier transform,'' Comput. Electron. Agric., vol. 135, pp. 115–127, Apr. 2017.
  \bibitem{Long_CNN} J. Long, E. Shelhamer and T. Darrell, ``Fully convolutional networks for semantic segmentation,'' 2015 IEEE Conference on Computer Vision and Pattern Recognition (CVPR), Boston, MA, USA, 2015, pp. 3431-3440, doi: 10.1109/CVPR.2015.7298965.
  \bibitem{Carson_Blobworld} C. Carson, S. Belongie, H. Greenspan and J. Malik, ``Blobworld: image segmentation using expectation-maximization and its application to image querying,'' in IEEE Transactions on Pattern Analysis and Machine Intelligence, vol. 24, no. 8, pp. 1026-1038, Aug. 2002, doi: 10.1109/TPAMI.2002.1023800.
  \bibitem{Zhang_MRF} Y. Zhang, M. Brady, and S. Smith, ``Segmentation of brain MR images through a hidden Markov random field model and the expectation-maximization algorithm,'' IEEE Trans. Med. Imag., vol. 20, no. 1, pp. 45–57, Jan. 2001.
  \bibitem{Zheng_OMRF} C. Zheng, X. Pan, X. Chen, X. Yang, X. Xin, and L. Su, ``An Object-Based Markov Random Field Model with Anisotropic Penalty for Semantic Segmentation of High Spatial Resolution Remote Sensing Imagery,'' Remote Sens., vol. 11, no. 23, Art. no. 2878, 2019.
  \bibitem{Yao_OMRF} H. Yao, X. Wang, L. Zhao, M. Tian, Z. Jian, L. Gong, and B. Li, ``An Object-Based Markov Random Field with Partition-Global Alternately Updated for Semantic Segmentation of High Spatial Resolution Remote Sensing Image,'' Remote Sens., vol. 14, no. 1, Art. no. 127, Jan. 2022, doi: 10.3390/rs14010127.
  \bibitem{ghosh_ddCRP} S. Ghosh, A. Ungureanu, E.Sudderth, and D. Blei, ``Spatial distance dependent Chinese restaurant processes for image segmentation,'' Advances in Neural Information Processing Systems, 24, 2011.
  \bibitem{wang_isita2024} Y. Wang, K. Shimada and T. Matsushima, ``An Efficient Image Segmentation Algorithm Based on Bayes Decision Theory and its Application for Region Detection,'' 2024 International Symposium on Information Theory and Its Applications (ISITA), Taipei, Taiwan, 2024, pp. 206-211.
  \bibitem{nakahara_quadtree} Y. Nakahara and T. Matsushima, ``A Stochastic Model for Block Segmentation of Images Based on the Quadtree and the Bayes Code for It,'' Entropy, vol.23, no.8, 2021.
  \bibitem{nakahara_full_rooted_trees} Y. Nakahara, S. Shota, A. Kamatsuka and T. Matsushima, ``Probability Distribution on Full Rooted Trees,'' Entropy, vol.24, no.3, 2022.
  \bibitem{Dryden_ShapePrior} I. L. Dryden and K. V. Mardia, ``Statistical Shape Analysis,'' Chichester, U.K.: Wiley, 1998.
  \bibitem{Gould_CRF} S. Gould, R. Fulton and D. Koller, ``Decomposing a scene into geometric and semantically consistent regions,'' 2009 IEEE 12th International Conference on Computer Vision, Kyoto, Japan, 2009, pp. 1-8.
  \bibitem{Oe_2DAR} S. Oe, ``Texture segmentation method by using two-dimensional AR model and Kullback information,'' Pattern Recognit., vol. 26, no. 2, pp. 237–244, Feb. 1993.
  \bibitem{berger_bayes_decision_theory} James O. Berger, ``Statistical Decision Theory and Bayesian Analysis,” Springer, 1985.
  \bibitem{Dempster_GMMEM} A. P. Dempster, N. M. Laird, and D. B. Rubin, ``Maximum likelihood from incomplete data via the EM algorithm,'' J. Roy. Statist. Soc. Ser. B (Methodol.), vol. 39, no. 1, pp. 1–22, 1977.
  \bibitem{IKONOS} National Aeronautics and Space Administration, ``IKONOS-2,'' NASA Open Data Portal, Sep. 20, 2019. [Online]. Available: https://data.nasa.gov/dataset/ikonos-2. Accessed: Apr. 30, 2026.
  \end{thebibliography}
\end{document}
